\documentclass[11pt]{article}
\usepackage{amsmath,amssymb,geometry,hyperref}
\geometry{margin=1in}
\newcommand{\Tr}{\operatorname{Tr}}
\begin{document}
\title{Analytical fourth-order AC states and the optimal subspace dimension at seventeen qubits}
\author{Catalogue calculation}
\date{30 October 2026}
\maketitle

Use normalized Dicke states with
$J_z|D_{17}^{(r)}\rangle=(17/2-r)|D_{17}^{(r)}\rangle$.
Set
\[
 q=\sqrt{20995},\qquad v=\sqrt{12780q-575935}>0,
\]
\[
 A=\frac{710-9q-iv}{756\sqrt{15}},\qquad
 B=\frac{2\sqrt{21}}{27},\qquad
 C=\frac{2(9q-2840-2iv)}{189\sqrt{4290}}.
\]
Then
\begin{align*}
 |\psi_+\rangle={}&A\bigl(|D_{17}^{(1)}\rangle+|D_{17}^{(16)}\rangle\bigr)\\
 &+B\bigl(|D_{17}^{(4)}\rangle+|D_{17}^{(13)}\rangle\bigr)\\
 &+C\bigl(|D_{17}^{(7)}\rangle+|D_{17}^{(10)}\rangle\bigr)
\end{align*}
is normalized and exactly fourth-order anticoherent. Its coefficientwise
complex conjugate $|\psi_-\rangle$ is a second such state. Each gives a
one-dimensional AC subspace; their span is not asserted to be AC.

This construction uses the same support pattern as the numerical seventeen-qubit
state in \texttt{ACLiege.jl}, together with an additional palindromic condition.
It is an independent analytical representative. Equality or rotation equivalence
to that numerical representative is not asserted.

\section{Derivation and exact verification of the state}
The probabilities are
\[
 |A|^2=\frac{277}{1458},\qquad |B|^2=\frac{28}{243},\qquad
 |C|^2=\frac{142}{729}.
\]
For palindromic support these are the unique probabilities matching the
zeroth, second and fourth diagonal spin moments of $I_{18}/18$.
Odd moments through order four vanish by reflection.

Only shifts divisible by three survive. Through order four the remaining
conditions are therefore
\[
 \langle J_+^3\rangle=0,\qquad
 \langle J_+^3(J_z+3/2)\rangle=0.
\]
Put $K_1=48\sqrt{35}$, $K_2=6\sqrt{10010}$, and
\[
 D=\frac{5680}{81},\qquad
 u=710-9q,\qquad x=\frac{8u}{243},\qquad y=\frac{8v}{243}.
\]
The coefficients above obey
\[
 K_1\overline A B=x+iy=:X,\qquad
 K_2\overline B C=-D-x-2iy=:Y,
\]
\[
 |X|^2=\frac{34746880}{19683},\qquad
 |Y|^2=\frac{159199040}{19683}.
\]
These identities follow from $q^2=20995$ and $u^2+v^2=1628760$.
The two remaining coherences become
\[
 2\Re(X+Y)+720|C|^2=0,\qquad
 12i\Im X+6i\Im Y=0.
\]
This proves the fourth-order conditions. The radicand is strictly positive
because $q>144$ and $12780\cdot144-575935>0$.

For a direct independent certificate, define the bipartite amplitude matrix
\[
 (M_t)_{ab}=c_{a+b}
 \sqrt{\frac{\binom ta\binom{17-t}b}{\binom{17}{a+b}}},\qquad
 0\le a\le t,\quad0\le b\le17-t.
\]
Exact substitution gives
\[
 M_4M_4^\dagger=\frac{I_5}{5}.
\]
The fifth reduction has diagonal entries $1/6$. Its only nonzero
upper-triangular off-diagonal entries are
\[
 (\rho_5)_{03}=(\rho_5)_{25}=-\frac{\sqrt{8398}}{9282},\qquad
 (\rho_5)_{14}=\frac{\sqrt{20995}}{9282}.
\]
Consequently
\[
 \Tr\rho_5^2=\frac{5443}{32487}
 =\frac16+\frac{19}{21658}>\frac16,
\]
so the order is exactly four. Complex conjugation preserves all these
reduced-state identities. The verifiers check every entry for both signs.

\section{The maximal fourth-order AC-subspace dimension is one}
The tensor-Gram bound gives $9k\le18$, so only $k=2$ needs exclusion.
Assume a rank-two fourth-order AC projector $P$ exists on spin $17/2$.
Equality in the tensor-Gram bound makes the images of its basis under
Hilbert--Schmidt-orthonormal spin tensors of ranks zero, one and two a
complete orthogonal basis. If $T_{\ell m}$ are those tensors, define the
covariant completely positive map
\[
 \mathcal E_2(A)=\sum_{\ell=0}^2\sum_{m=-\ell}^{\ell}
 T_{\ell m}A T_{\ell m}^\dagger.
\]
The normalization $\Tr(T_{\ell m}^\dagger T_{\ell'm'})=
\delta_{\ell\ell'}\delta_{mm'}$ gives
$\mathcal E_2(P)=I_{18}/18$.

On tensor rank $L$, put $x=L(L+1)$. The channel multiplier is
\[
 e_{2,L}=\frac{3\bigl(10x^2-3860x+310080\bigr)}{18\cdot323\cdot320}
 =\frac{30(x-272)(x-114)}{18\cdot323\cdot320}.
\]
Among integer ranks $1\le L\le17$, it vanishes only at $L=16$.
Therefore
\[
 P=\frac{I_{18}}9+X,\qquad X\text{ has only tensor rank }16.
\]

Let $\mathcal R_1$ trace out one physical qubit. The first-order AC cross
conditions imply that
\[
 Q=2\mathcal R_1(P)
\]
is a rank-four orthogonal projector on the seventeen-dimensional symmetric
sixteen-qubit space: the two Schmidt partners of each of the two code
basis vectors form four orthogonal unit vectors. Rotation covariance gives
\[
 Q=\frac4{17}I_{17}+Y,\qquad Y\text{ has only tensor rank }16.
\]
Here $Y$ is Hermitian and nonzero. From $Q^2=Q$,
\[
 Y^2=\frac9{17}Y+\frac{52}{289}I_{17},\qquad
 \Tr Y^2=\frac{52}{17}.
\]

Let $F$ be the unitary linear part of spin-eight time reversal.
Time-reversed vectorization identifies a rank-sixteen operator with the
maximal-spin subspace of $V_8\otimes V_8$. Consequently
\[
 |\Omega\rangle=\frac{\operatorname{vec}(YF)}{\sqrt{\Tr Y^2}}
\]
is a pure symmetric thirty-two-qubit state. Its sixteen-qubit reduction is
\[
 \frac{Y^2}{\Tr Y^2}=\frac{I_{17}}{17}+\frac9{52}Y.
\]
Tracing one more qubit eliminates tensor rank sixteen, since the output
operator space has ranks at most fifteen. Thus its fifteen-qubit reduction
is $I_{16}/16$: $|\Omega\rangle$ would be fifteenth-order AC.

The ordinary tensor-Gram bound for a pure fifteenth-order AC state requires
\[
 \sum_{\ell=0}^7(2\ell+1)=64\le32+1=33,
\]
which is impossible. Hence $k=2$ does not exist. The radical state already
supplies $k=1$, proving
\[
 \boxed{K(17,4)=1.}
\]
This upper-bound argument applies to arbitrary complex subspaces, without
assuming palindromicity, sparse support or any constellation symmetry.

\section{Rotations and conjugate representatives}
The sparse support gives a threefold axis, and palindromicity gives a
perpendicular half-turn. Thus the proper rotational symmetry contains $D_3$.
The nonzero fifth-rank multipole, implied by fourth-order AC and failure of
fifth-order AC, rules out tetrahedral symmetry: its rank-five invariant
multiplicity is $(11-8-3)/12=0$. Octahedral and icosahedral groups contain a
tetrahedral subgroup and are excluded as well. The support spacing permits
no axial rotation of order greater than three. Since the constellation is
not confined to an axis, continuous rotation symmetry is impossible.
The full proper group is therefore $D_3$, with a unique threefold axis.

Any rotation relating the two conjugate representatives must preserve this
axis up to sign. The quantity $c_1c_{16}/(c_4c_{13})=A^2/B^2$ is invariant
under these rotations and a global phase, but has nonzero imaginary part.
Conjugation changes it to its distinct conjugate. The two representatives
therefore give inequivalent proper-rotation orbits. No classification of
all fourth-order states at seventeen qubits is claimed.

\section{Exact Majorana multiplicities}
Use $P(z)=\sum_r(-1)^r\sqrt{\binom{17}r}\,c_rz^r$.
With $u=710-9q$ and $w=iv$, the polynomial factors as
\[
 P(z)=-\sqrt{17}A\,z\,Q(z^3),
\]
\[
 Q(y)=1+Ly+My^2-My^3-Ly^4-y^5,
\]
\[
 L=-\frac{2(u+w)}{277},\qquad
 M=-\frac{16(2130+u+2w)(u+w)}{1628760}.
\]
In particular,
\[
 Q(y)=(1-y)\bigl[1+(L+1)y+(M+L+1)y^2+(L+1)y^3+y^4\bigr].
\]
The exact field $\mathbb Q(w)$ contains $q$ because
$q=(575935-w^2)/12780$.
Exact polynomial arithmetic in this quartic field verifies
$\gcd(Q,Q')=1$ and $Q(0)=1$. Thus the constellation has seventeen
distinct positions: one simple root at zero, fifteen nonzero simple roots,
and one simple projective root at infinity. Replacing $w$ by $-w$ proves
the same statement for the conjugate representative. The geometry verifier
checks the exact Dicke-to-polynomial identities before evaluating roots;
its displayed sphere coordinates are only a numerical rendering of these
certified simple roots.
\end{document}
